%----------------------------------------
\begin{frame}
  \frametitle{Part 5}

\icontp[3.2]

\vspace{0.45cm}

{\Large\bf Theorem proving: all executions, no bound}

\vspace{0.4cm}

{\small \ours{\fcheck{}} using \lean{} 4, \ours{\tasksat{}} using \zthree{},\\[0.2cm]
\ours{\esbmcpy{}} using \esbmc{}}

\end{frame}

%----------------------------------------
\begin{frame}[fragile]
  \frametitle{\fcheck{}: proving an \fprime{} application} \small

\newtag{} \ {\bf \fcheck{}} --- a {\em new} tool of ours, using \lean{} 4. The
hysteresis logic of {\tt HeaterController}, in flight C++:

\begin{lstlisting}[style=cpp]
if (m_controlMode != ControlMode::AUTOMATIC) return;
HeaterState newState = m_heaterState;
if      (avgTemp < (m_targetTemp - m_hysteresis)) newState = ON;
else if (avgTemp > (m_targetTemp + m_hysteresis)) newState = OFF;
if (newState != m_heaterState) commandHeater(newState);
\end{lstlisting}

\vspace{-0.15cm}

and the \lean{} model --- \iconllm[0.85]\ {\bf LLM-generated} --- tracking it line
for line:

\begin{lstlisting}[language=lean4]
def avgTempIn_handler (s : ControllerState) (d : AvgTempData) :=
  if s.controlMode != ControlMode.AUTOMATIC then s
  else let newState :=
    if d.avgTemp < s.targetTemp - s.hysteresis then ON
    else if d.avgTemp > s.targetTemp + s.hysteresis then OFF
    else s.heaterState
  if newState != s.heaterState then commandHeater s newState else s
\end{lstlisting}

\end{frame}

%----------------------------------------
\begin{frame}[fragile]
  \frametitle{\fcheck{}: what the proof attempt found} \small

\begin{lstlisting}[language=lean4]
theorem cold_reading_turns_heater_on (s) (d)
    (ha    : s.controlMode = ControlMode.AUTOMATIC)
    (hcold : d.avgTemp < s.targetTemp - s.hysteresis) :
    (avgTempIn_handler s d).heaterState = HeaterState.ON
\end{lstlisting}

Too cold $\Rightarrow$ heater ON, for {\em every} state and {\em every} reading.

\pause
\vspace{0.4cm}

\begin{block}{Arbitrary operator behaviour}\small
\begin{lstlisting}[language=lean4]
theorem executeTrace_preserves_invariant
    (sys) (trace : List SystemEvent)
    (hInv : SystemInvariant sys) :
    SystemInvariant (executeTrace sys trace)
\end{lstlisting}
{\scriptsize No sequence of commands and ticks --- of {\em any} length --- can break the
invariant. \spin{} would have had to bound that list.}
\end{block}

\end{frame}

%----------------------------------------
\begin{frame}[fragile]
  \frametitle{\tasksat{}: proof without a proof script} \small

\newtag{} \ a {\em new} tool for {\bf planning and scheduling}, inspired by MEXEC.
A {\bf task net} declares timelines --- states, resources, rates --- and tasks with
windows, preconditions and impacts:

\begin{lstlisting}[language=tasknetdsl,frame=single]
  battery : rate [10.0, 100.0] bounds [0.0, 100.0] = 60.0;

  task drive {
    startrange [100, 120];  durationrange [30, 40];
    pre     { battery in [60.0, 100.0]; }
    impacts { maint { battery +~ -1.5; }  post { location = target; } }
  }

  constraints { always not (active(charge) and active(drive)); }
  properties  { eventually (location = target and data >= 30); }
\end{lstlisting}

\pause

\begin{block}{Two problems, one encoding}\small
\[
\text{schedule:}\quad \exists\, \mathit{init}.\ \exists\, \mathit{sched}.\ C(\mathit{init},\mathit{sched})
\qquad
\text{verify:}\quad \forall\, \mathit{init}.\ \forall\, \mathit{sched}.\ C \rightarrow \varphi
\]
\end{block}

Flip one quantifier and the planner becomes a {\bf prover} --- discharged
symbolically, for {\em every} initial battery level and start time, not a state graph.

\end{frame}

%----------------------------------------
\begin{frame}
  \frametitle{\esbmcpy{}: verifying Python} \small

\newtag{} \ Why Python: we are building a {\bf Python model of an autonomous system}
in the {\bf PAAS} project. And Python has no model checker of the kind \spin{} is ---
so {\bf translate it into something that does}, and let \esbmc{} decide.

\vspace{0.35cm}

{\centering
\begin{tikzpicture}[
  font=\scriptsize, >=stealth,
  b/.style={draw, rounded corners=1.5pt, align=center, inner sep=4pt, fill=black!4},
  lab/.style={font=\tiny, inner sep=2pt}
]
  \node[b] (py)  at (0,0)    {Python};
  \node[b] (c)   at (3.5,0)  {verifiable\\C / C++};
  \node[b, fill=iconcol!10] (e) at (7.0,0) {\esbmc{}};
  \node[b, fill=green!14] (ok)  at (9.9,0.55)  {ok};
  \node[b, fill=red!10]   (bad) at (9.9,-0.55) {counterexample};

  \draw[->] (py) -- node[lab, above] {\iconllm[0.6]\ LLM} (c);
  \draw[->] (c)  -- (e);
  \draw[->] (e)  -- (ok);
  \draw[->] (e)  -- (bad);
\end{tikzpicture}\par}

\vspace{0.35cm}

Three routes to the C: {\bf Shedskin} for what converts statically, an {\bf LLM} for
what does not, and {\bf dynamic tracing} for what neither can reach.

\vspace{0.2cm}

Driving it: \ours{\eva{}}, the {\bf Enhanced Verification Agent} --- which also runs
the static analysers and the runtime checks, and iterates.

\vspace{0.2cm}

\begin{block}{}\small
\centering
The LLM does the translating. \esbmc{} does the deciding.
\end{block}

\vspace{0.25cm}

{\scriptsize Largely developed by GenAI and the University of Manchester.}

\end{frame}
